\documentclass[tikz,border=8pt]{standalone}
\usepackage{tikz}
\usepackage{amsmath}
\usepackage{amssymb}
\usepackage{pgfplots}
\pgfplotsset{compat=1.18}
\usepackage{ctex}
\usetikzlibrary{calc,positioning,arrows.meta}

% LC 76 Minimum Window Substring — 固定/可变滑窗
% s = "ADOBECODEBANC"，t = "ABC"
% 3 个快照：扩张找合法窗 → 收缩到最小 → 继续扩张

\begin{document}
\begin{tikzpicture}

\def\cs{0.85}  % 格子步长

%% ── 辅助：画格子 ─────────────────────────────────────────
% \charcell{x}{y}{char}{fillcolor}
\newcommand{\charcell}[4]{%
  \node[draw, minimum size=0.8cm, inner sep=0pt,
        fill=#4, font=\small\bfseries] at (#1,#2) {#3};
}

%% ── 标题 ─────────────────────────────────────────────────
\node[font=\large\bfseries] at (5.0, 0.8)
  {LC 76：固定/可变滑窗（Minimum Window Substring）};
\node[font=\small, gray] at (5.0, 0.3)
  {s = ``ADOBECODEBANC''，t = ``ABC''};

%% 字符串：A D O B E C O D E B A N C
%%         0 1 2 3 4 5 6 7 8 9 10 11 12

%% ── 辅助宏：一行画 s 的格子 ─────────────────────────────
%% 颜色按列表传入，共 13 格
%% 手动列出各快照（避免复杂参数）

%% ═════════════════════════════════════════════════════════
%% 快照 0：扩张 r 直到窗 [0..5] = "ADOBEC" 覆盖 t="ABC"
%% ═════════════════════════════════════════════════════════
\begin{scope}[yshift=-0.5cm]
  % 格子颜色：0=A 绿覆盖，3=B 绿覆盖，5=C 绿覆盖，其余窗内淡绿
  \charcell{0*\cs}{0}{A}{green!25}   % 0
  \charcell{1*\cs}{0}{D}{green!10}   % 1
  \charcell{2*\cs}{0}{O}{green!10}   % 2
  \charcell{3*\cs}{0}{B}{green!25}   % 3
  \charcell{4*\cs}{0}{E}{green!10}   % 4
  \charcell{5*\cs}{0}{C}{green!25}   % 5
  \charcell{6*\cs}{0}{O}{white}      % 6
  \charcell{7*\cs}{0}{D}{white}      % 7
  \charcell{8*\cs}{0}{E}{white}      % 8
  \charcell{9*\cs}{0}{B}{white}      % 9
  \charcell{10*\cs}{0}{A}{white}     % 10
  \charcell{11*\cs}{0}{N}{white}     % 11
  \charcell{12*\cs}{0}{C}{white}     % 12
  % 下标
  \foreach \i in {0,...,12}{
    \node[font=\tiny,gray] at (\i*\cs,-0.55) {\i};
  }
  % 绿色窗框
  \draw[green!50!black,thick,rounded corners=2pt]
    (-0.43,-0.43) rectangle (5*\cs+0.43,0.43);
  % l 和 r 指针
  \draw[->,>=Stealth,green!60!black,thick] (0,-0.9) -- (0,-0.5);
  \node[font=\scriptsize,green!60!black,below] at (0,-0.95) {l=0};
  \draw[->,>=Stealth,green!60!black,thick] (5*\cs,-0.9) -- (5*\cs,-0.5);
  \node[font=\scriptsize,green!60!black,below] at (5*\cs,-0.95) {r=5};
  % 目标字符标注
  \foreach \x/\c in {0/A,3/B,5/C}{
    \node[font=\tiny,green!60!black] at (\x*\cs,0.6) {$\checkmark$};
  }
  % 注释框
  \node[draw=gray!50,fill=white,thin,rounded corners=2pt,
        font=\scriptsize,inner sep=4pt,align=left]
    at (14.8,-0.35)
    {扩张 r 至 5，窗 ``ADOBEC''\\
     已包含 A、B、C（t 全覆盖）\\
     当前窗长 = \textbf{6}};
\end{scope}

%% ═════════════════════════════════════════════════════════
%% 快照 1：收缩 l，最小合法窗 [3..5] = "BEC"? 不对，
%% 实际收缩：l=0→"DOBEC"? 检查：A已不在…
%% 正确最小窗：收缩到 l 使得 A 消失时停止，即 l=0(A需保留)
%% 等等，窗 [0..5] = "ADOBEC"（6字符），收缩 l：
%%   l=1: "DOBEC"——缺 A，不合法，停止，最小窗 = "ADOBEC"（l=0）
%% 然后记录 minWin = "ADOBEC"(len=6), l++: l=1（窗不合法），继续扩 r
%% 实际：用 l=0..5 窗 → 收缩后最小 = [0,5] len=6
%% 快照1展示：已记录 best，l 移到 1，窗不合法（灰色）
%% ═════════════════════════════════════════════════════════
\begin{scope}[yshift=-3.1cm]
  \charcell{0*\cs}{0}{A}{gray!20}   % 已移出（l 跳过）
  \charcell{1*\cs}{0}{D}{gray!10}   % 收缩起点
  \charcell{2*\cs}{0}{O}{gray!10}
  \charcell{3*\cs}{0}{B}{gray!15}
  \charcell{4*\cs}{0}{E}{gray!10}
  \charcell{5*\cs}{0}{C}{gray!15}
  \charcell{6*\cs}{0}{O}{white}
  \charcell{7*\cs}{0}{D}{white}
  \charcell{8*\cs}{0}{E}{white}
  \charcell{9*\cs}{0}{B}{white}
  \charcell{10*\cs}{0}{A}{white}
  \charcell{11*\cs}{0}{N}{white}
  \charcell{12*\cs}{0}{C}{white}
  \foreach \i in {0,...,12}{
    \node[font=\tiny,gray] at (\i*\cs,-0.55) {\i};
  }
  % 已记录最优（虚线框）
  \draw[green!50!black,dashed,thick,rounded corners=2pt]
    (-0.43,-0.43) rectangle (5*\cs+0.43,0.43);
  \node[font=\tiny,green!50!black] at (2.5*\cs,0.62) {已记录 best = ``ADOBEC''(len=6)};
  % 当前 l=1，收缩后缺 A，窗不合法（红框）
  \draw[red!50!black,thick,rounded corners=2pt]
    (\cs-0.43,-0.43) rectangle (5*\cs+0.43,0.43);
  \draw[->,>=Stealth,orange!80!black,thick] (\cs,-0.9) -- (\cs,-0.5);
  \node[font=\scriptsize,orange!80!black,below] at (\cs,-0.95) {l=1};
  \draw[->,>=Stealth,green!60!black,thick] (5*\cs,-0.9) -- (5*\cs,-0.5);
  \node[font=\scriptsize,green!60!black,below] at (5*\cs,-0.95) {r=5};
  \node[draw=gray!50,fill=white,thin,rounded corners=2pt,
        font=\scriptsize,inner sep=4pt,align=left]
    at (14.8,-0.35)
    {收缩：l++ $\Rightarrow$ l=1\\
     窗 ``DOBEC'' 缺 A，不合法\\
     $\Rightarrow$ 停止收缩，继续扩 r};
\end{scope}

%% ═════════════════════════════════════════════════════════
%% 快照 2：继续扩 r 到 10（A 重新进入），收缩到最小 [9,12]
%% 窗 [9..12] = "BANC"，len=4，更新 best
%% ═════════════════════════════════════════════════════════
\begin{scope}[yshift=-5.7cm]
  \charcell{0*\cs}{0}{A}{gray!15}
  \charcell{1*\cs}{0}{D}{gray!15}
  \charcell{2*\cs}{0}{O}{gray!15}
  \charcell{3*\cs}{0}{B}{gray!15}
  \charcell{4*\cs}{0}{E}{gray!15}
  \charcell{5*\cs}{0}{C}{gray!15}
  \charcell{6*\cs}{0}{O}{gray!15}
  \charcell{7*\cs}{0}{D}{gray!15}
  \charcell{8*\cs}{0}{E}{gray!15}
  \charcell{9*\cs}{0}{B}{green!25}   % 9
  \charcell{10*\cs}{0}{A}{green!25}  % 10
  \charcell{11*\cs}{0}{N}{green!10}  % 11
  \charcell{12*\cs}{0}{C}{green!25}  % 12
  \foreach \i in {0,...,12}{
    \node[font=\tiny,gray] at (\i*\cs,-0.55) {\i};
  }
  % 绿色窗框 [9..12]
  \draw[green!50!black,thick,rounded corners=2pt]
    (9*\cs-0.43,-0.43) rectangle (12*\cs+0.43,0.43);
  % 目标字符标注
  \foreach \x in {9,10,12}{
    \node[font=\tiny,green!60!black] at (\x*\cs,0.6) {$\checkmark$};
  }
  % l 和 r 指针
  \draw[->,>=Stealth,green!60!black,thick] (9*\cs,-0.9) -- (9*\cs,-0.5);
  \node[font=\scriptsize,green!60!black,below] at (9*\cs,-0.95) {l=9};
  \draw[->,>=Stealth,green!60!black,thick] (12*\cs,-0.9) -- (12*\cs,-0.5);
  \node[font=\scriptsize,green!60!black,below] at (12*\cs,-0.95) {r=12};
  % 注释框
  \node[draw=gray!50,fill=white,thin,rounded corners=2pt,
        font=\scriptsize,inner sep=4pt,align=left]
    at (14.8,-0.35)
    {扩张至 r=12，收缩至 l=9\\
     窗 ``BANC''（含 A、B、C）\\
     len = \textbf{4} $<$ 6，更新 best\\
     \textbf{最终答案}：``BANC''};
\end{scope}

%% ── 图例 ─────────────────────────────────────────────────
\begin{scope}[yshift=-8.1cm]
  \node[draw,minimum size=0.5cm,fill=green!25,inner sep=0pt] at (0.0,0) {};
  \node[font=\scriptsize,right] at (0.35,0) {包含 t 目标字符};
  \node[draw,minimum size=0.5cm,fill=green!10,inner sep=0pt] at (3.2,0) {};
  \node[font=\scriptsize,right] at (3.55,0) {窗内（非目标）};
  \node[draw,minimum size=0.5cm,fill=gray!15,inner sep=0pt] at (6.2,0) {};
  \node[font=\scriptsize,right] at (6.55,0) {窗外/已排除};
\end{scope}

\end{tikzpicture}
\end{document}
